% Source-aligned English for printed p.127.
\SourcePage{127}

\LayerTarget{tohoku:section:1.4}
\noindent\textbf{1.4. Abelian categories.}
An \textit{abelian category} is an additive category $\mathbf{C}$ satisfying the two
following additional axioms (which are self-dual):

\LayerTarget{tohoku:section:1.4:AB1}
\noindent AB 1) \textit{Every morphism admits a kernel and a cokernel}
(cf. 1.3.)

\LayerTarget{tohoku:section:1.4:AB2}
\noindent AB 2) \textit{Let $u$ be a morphism in $\mathbf{C}$. Then the
canonical morphism
$\bar u:\mathrm{Coim}\,u\to\mathrm{Im}\,u$ \textup{(cf. 1.3)} is an
isomorphism.}

\LayerTarget{tohoku:section:1.4:p1}
It follows in particular that \textit{a bijective morphism is an isomorphism}.
Note that there are many additive categories satisfying AB 1) in which the
morphisms $\bar u:\mathrm{Coim}\,u\to\mathrm{Im}\,u$ are always
bijective without necessarily being isomorphisms. This is the case, for
example, for the additive category of separated topological modules over a
given topological ring, with continuous homomorphisms as morphisms, or for the
category of \textit{filtered} abelian groups. Another, less obvious example is
the additive category of holomorphic fibre spaces with vector-space fibre over
a holomorphic manifold of complex dimension 1. These are therefore
nonabelian additive categories.

\LayerTarget{tohoku:section:1.4:p2}
If $\mathbf{C}$ is an abelian category, the entire usual formalism of diagrams of
homomorphisms between abelian groups can be developed anew by replacing the
homomorphisms with morphisms in $\mathbf{C}$, provided that one has in view only
properties ``of finite character'', i.e. properties not involving infinite
direct sums or direct products (for which special precautions are necessary;
see No.~5). We shall confine ourselves here to indicating a few particularly
important facts, referring to [3] for further details.

\LayerTarget{tohoku:section:1.4:p3}
In what follows, we place ourselves in a fixed abelian category $\mathbf{C}$. Let
$A\in\mathbf{C}$. To every subobject of $A$ associate its cokernel (which is therefore
a quotient of $A$), and to every quotient object of $A$ associate its kernel
(which is therefore a subobject of $A$). This gives \textit{a one-to-one
correspondence between the class of subobjects of $A$ and the class of quotient
objects of $A$}; moreover, this correspondence is an \textit{order
anti-isomorphism} for the natural order relations. Furthermore, the subobjects
of $A$ form a class with a \textit{lattice} structure (and hence so do the quotient objects): if $P$
and $Q$ are two subobjects of $A$, their supremum is the image of the direct
sum $P+Q$ under the morphism whose components are the canonical injections of
$P$ and $Q$ into $A$, and their infimum is the kernel of the morphism from $A$
to the product $(A/P)\times(A/Q)$ whose components are the canonical surjections from $A$ onto $A/P$
and $A/Q$. (In accordance with usage, $A/P$ denotes the quotient of $A$
corresponding to the subobject $P$; a dual notation such as $A\backslash R$ for
the subobject of $A$ corresponding to a quotient object $R$ seems natural.)
There are dual interpretations for the infimum or the supremum of two quotient
objects of $A$.

\LayerTarget{tohoku:section:1.4:p4}
Let $u:A\to B$ be a morphism. If $A'$ is a subobject of $A$, the
image of $A'$ under $u$, denoted $u(A')$, is defined as $\mathrm{Im}\,ui$,
where $i$ is the canonical injection $A'\to A$. Dually, one
defines the inverse image $u^{-1}(B')$ of a quotient $B'$ of $B$; it is a
quotient of $A$. If $B'$ is now a subobject of $B$, the inverse image of $B'$
under $u$, denoted $u^{-1}(B')$, is defined as the kernel of $ju$, where $j$
is

\TohokuContinuation{128}
